\documentclass[aspectratio=169,usenames,dvipsnames]{beamer}
\beamertemplatenavigationsymbolsempty
% \documentclass{beamer}

% UNCOMMENT FOR HANDOUTS
% \usepackage{pgfpages}
% \pgfpagesuselayout{2 on 1}[landscape,letterpaper,border shrink=5mm]
% \setbeameroption{show notes}

\mode<presentation>
\usetheme{Boadilla}
\usecolortheme{spruce}
\usepackage{../LizStyle,ulem}
\usepackage{color}
\usepackage[color,matrix,arrow]{xy}
% \input xy
% \xyoption{all}
\usepackage{tikz}
\usepackage{tikz-cd}
\tikzset{commutative diagrams/diagrams={ampersand replacement=\&}}

\AtBeginSection{\frame{\sectionpage}}


% ==========Command for putting comments and removing them=======
% Visible:
\newcommand{\InstructorNotes}[1]{\textit{\color{orange}#1}}
\newcommand{\InstructorList}[1]{
\begin{itemize}
	\color{orange}
	#1
\end{itemize}
}
% Removed:
% \renewcommand{\InstructorNotes}[1]{}
% \renewcommand{\InstructorList}[1]{}


\newcommand{\bvfill}{\vskip0pt plus 1filll}

\usepackage{multimedia}
\usepackage{graphbox}

% \usepackage{algpseudocode}

\graphicspath{ {Figures/} {../Figures/}{../../Figures/} }



\begin{document}
%\pdfoptionpdfminorversion 5
%\logo{\includegraphics[height=1.5cm]{dukeshield.pdf}}
\title[]{Ch 10.1: Neural Nets}
\subtitle{Lecture 29 - CMSE 381}
\date{Mon, Nov 18, 2024}



\institute[MSU-CMSE]{Michigan State University\\
			::\\
          Dept of Computational Mathematics, Science \& Engineering}
\author[Dr.~Munch]{Prof.~Elizabeth Munch}

\begin{frame}
 \titlepage


\end{frame}
%--------------------------
\begin{frame}[c]{Announcements }

\begin{columns}
\column{.45\textwidth}
\centering

	\textbf{Last time:}
	\begin{itemize}
	\item SVM 
	\end{itemize}

	\vspace{.2in}

	\textbf{This lecture:}
	\begin{itemize}
		\item Feed Forward Neural Nets
\end{itemize}

\textbf{Announcements:}
	\begin{itemize}
	\item Homework \#9 due Sunday, Nov 24 
	\item No class wednesday Nov 27. Virtual office hours by request.

	\end{itemize}

\column{.4\textwidth}
\centering
\includegraphics[width = \textwidth]{Day01-Schedule-Test3.png}

\end{columns}

\end{frame}
%--- Next Frame ---%


%==========
\section{Neural Nets}
%==========

\begin{frame}[t]{The idea}

\begin{columns}
\column{.45\textwidth}
\centering
\InstructorList{
  \item Draw input of vector $p$ variables $X = (X_1,\cdots,X_p)$
	\item Draw black box nonlinear function $f(X)$
	\item Spit out response prediction $Y$
	\\
	\item We've had nonlinear stuff before like trees, boosting, GANs
	\item Difference here is the structure
}

\column{.45\textwidth}
\centering

\end{columns}

\end{frame}
%--- Next Frame ---%
\begin{frame}[t]{Feed Forward Neural Network: The cartoon}

\begin{columns}
\column{.45\textwidth}
\centering
\includegraphics[width = \textwidth, align = c]{Fig10-1}
\column{.45\textwidth}
\centering
\InstructorList{
  \item Input layer is starting data
	\item Each arrow means taking combo of those values with weights to get next value
	\item Here there is one hidden layer then the output
	\item Get to pick how many hidden units $K$, here we have $K=5$
}

\end{columns}

\end{frame}
%--- Next Frame ---%
\begin{frame}[t]{A very simple example}
	{Computing $A_k$ for $(1,0)$}

\includegraphics[width = .4\textwidth, align = c]{Eq10-2}

\begin{columns}
\column{.45\textwidth}
\centering
\includegraphics[width = \textwidth, align = c]{DL-simple-ex}
\column{.45\textwidth}
\centering
\InstructorNotes{
\includegraphics[width = .5\textwidth, align = c]{DL-simple-ex-Input10}
\includegraphics[width = .5\textwidth, align = c]{DL-simple-ex-Input10-Soln}
}
\InstructorList{
  \item $A_k$ called activations
	\item Pick an activation function $g(z)$
}

\end{columns}

\end{frame}
%--- Next Frame ---%


\begin{frame}[t]{A very simple example}
	{Computing $Y$ for $(1,0)$}


$\displaystyle{f(X) = \beta_0 + \sum_{k=1}^K \beta_k A_k}$

\begin{columns}
\column{.45\textwidth}
\centering
\includegraphics[width = \textwidth, align = c]{DL-simple-ex}
\column{.45\textwidth}
\centering
\InstructorNotes{
\includegraphics[width = .5\textwidth, align = c]{DL-simple-ex-Input10}
\includegraphics[width = .5\textwidth, align = c]{DL-simple-ex-Input10-Soln}
}
\end{columns}

\end{frame}
%--- Next Frame ---%


\begin{frame}[t]{A very simple example}
	{Computing $Y$ for $(0,1)$}


\includegraphics[width = .4\textwidth, align = c]{Eq10-2}
\qquad
$\displaystyle{f(X) = \beta_0 + \sum_{k=1}^K \beta_k A_k}$

\begin{columns}
\column{.45\textwidth}
\centering
\includegraphics[width = \textwidth, align = c]{DL-simple-ex}
\column{.45\textwidth}
\centering
\InstructorNotes{
\includegraphics[width = .5\textwidth, align = c]{DL-simple-ex-Input01}
\includegraphics[width = .5\textwidth, align = c]{DL-simple-ex-Input01-Soln}
}
\end{columns}

\end{frame}
%--- Next Frame ---%
{
\setbeamercolor{background canvas}{bg=Apricot!25}
\begin{frame}[t]{A different example}
\begin{columns}
\column{.55\textwidth}
\centering
\begin{itemize}
	\item Draw the diagram for a neural net with input data points with $p=3$ (i.e., $(X_1,X_2,X_3)$) and two  units in the hidden layer. 
	\item Using the $\omega$ and $\beta$ matrices, what is the output predicted $Y$ for the point $(2,0, 1)$?
	\begin{equation*}
		\omega = 
		\begin{pmatrix}
			1 & 0 & -2  & 2\\
			-3 & 1 & 0 & -1
		\end{pmatrix}
		\qquad 
		\beta = 
		\begin{pmatrix}
			-1 \\
			-2 \\
			1
		\end{pmatrix}
	\end{equation*}
	\item Use the activation function 
\begin{equation*}
	g(z) = (z)_+ =
	\begin{cases}
		0 & \text{if }z<0 \\
		z & \text{else.}
	\end{cases}
\end{equation*}
	
\end{itemize}
\column{.35\textwidth}
\centering
\end{columns}
\end{frame}
}
%--- Next Frame ---%
{
\setbeamercolor{background canvas}{bg=Apricot!25}
\begin{frame}[t]{Extra space}
\begin{columns}
\column{.45\textwidth}
\centering
\InstructorNotes{
	\includegraphics[align = c, width = \textwidth]{NN-ExampleArchitecture.jpeg}
	Oops, those three shouldn't be labeled as A's, it should be X's.
	Also check the right, I did the relu computation wrong so some error propagated!		
}
\column{.45\textwidth}
\InstructorNotes{
	\includegraphics[align = c, width = \textwidth]{NN-ExampleCalc.jpeg}
}
\centering
\end{columns}
\end{frame}
}
%--- Next Frame ---%
\begin{frame}[t]{What if there's no activation function?}
\begin{columns}
\column{.45\textwidth}
\centering

\includegraphics[align = c, width = \textwidth]{NN-Example2-Architecture.jpeg}
\begin{equation*}
	\omega = 
	\begin{pmatrix}
		1 & 2 & 3  \\
		4 & 5 & 6 
	\end{pmatrix}
	\qquad 
	\beta = 
	\begin{pmatrix}
		-1 \\
		-2 \\
		1
	\end{pmatrix}
\end{equation*}
\column{.45\textwidth}
\centering
\InstructorNotes{
	Calculate with no activation function to show that the result is just some other linear function 
% 
	\includegraphics[align = c, width = \textwidth]{NN-Example2-Calc.jpeg}
}
\end{columns}
\end{frame}
%--- Next Frame ---%

\begin{frame}[t]{Choices for activation function }

	\begin{columns}
	\column{.45\textwidth}
	\centering
	Sigmoid:
	\begin{equation*}
		g(z) = \frac{e^z}{1+e^z} = \frac{1}{1+e^{-z}}
	\end{equation*}
	% \InstructorList{
	%   \item Converts linear function in probabilities in [0,1]
	% }

	\includegraphics[align = c, width = \textwidth]{Sigmoid-pytorch}
	
	
	\column{.45\textwidth}
	\centering
	ReLU: Rectified linear unit
	\begin{equation*}
		g(z) = (z)_+ =
		\begin{cases}
			0 & \text{if }z<0 \\
			z & \text{else.}
		\end{cases}
	\end{equation*}
	\includegraphics[align = c, width = \textwidth]{ReLU-pytorch}
	\end{columns}
	
	\end{frame}

\begin{frame}[t]{Matrix version: First layer}
	\begin{columns}
	\column{.25\textwidth}
	\centering
\includegraphics[width = \textwidth, align = c]{DL-simple-ex}
	
	\column{.65\textwidth}
	\centering
	\includegraphics[width = .6\textwidth, align = c]{Eq10-2}
\begin{equation*}
	A = g( \mathbf{W}\cdot \mathbf{X})
	\qquad 
	\mathbf{X}^T = (1 \; X_1\; X_2\; \cdots \; X_p)
\end{equation*}
	\begin{equation*}
	\end{equation*}
	\InstructorNotes{Caclulate out the matrix multiplication using the matrices at left to show that the equations are the same.}
	\end{columns}
	\end{frame}
%-----------end frame------------

\begin{frame}[t]{Matrix version: Output}
	\begin{columns}
	\column{.25\textwidth}
	\centering
\includegraphics[width = \textwidth, align = c]{DL-simple-ex}
	
	\column{.65\textwidth}
	\centering
$\displaystyle{f(X) = \beta_0 + \sum_{k=1}^K \beta_k A_k}$

\begin{equation*}
	Y = \beta \cdot \textbf{A}
	\qquad 
	\mathbf{A}^T = (1 \; A_1\; A_2\; \cdots \; A_K)
\end{equation*}
	\begin{equation*}
	\end{equation*}
	\end{columns}
	\end{frame}
%-----------end frame------------

\begin{frame}[t]{Now what?}
Choose parameters by minimizing RSS, $\sum_{i=1}^n(y_i - f(x_i))^2$
\begin{columns}
\column{.45\textwidth}
\centering
\textbf{Chosen in advance:}
\InstructorList{
  \item Number of layers (more on that next lecture)
	\item Number of hidden units
	\item Activation function $g(z)$
}


\column{.45\textwidth}
\centering
\textbf{Tuned by the model:}
\InstructorList{
  \item $\omega$
	\item $\beta$
}

\end{columns}

\end{frame}
%--- Next Frame ---%

{
\setbeamercolor{background canvas}{bg=Apricot!25}
\begin{frame}[t]{Coding}
\begin{columns}
\column{.45\textwidth}
\centering
\column{.45\textwidth}
\centering
\end{columns}
\end{frame}
}
%--- Next Frame ---%

% %==========
% \section{Multilayer Neural Networks}
% %==========

% \begin{frame}[t]{Multiple layers}
% \begin{columns}
% \column{.45\textwidth}
% \centering
% \includegraphics[width = \textwidth, align = c]{Fig10-4}

% \column{.45\textwidth}
% \centering
% \InstructorList{
%   \item Include more layers 
%   \item Can pick number of units per layer 
%   \item Each layer is linear combinations of previous
% }

% \end{columns}
% \end{frame}
% %--- Next Frame ---%

% \begin{frame}[t]{Hidden layers}

% \begin{columns}
% \column{.45\textwidth}
% \centering
% \includegraphics[width = \textwidth, align = c]{Fig10-4}
% \column{.45\textwidth}
% \centering
% \includegraphics[width = \textwidth, align = c]{Eq10-10}
% \includegraphics[width = \textwidth, align = c]{Eq10-11}
% \InstructorList{
%   \item Two hidden layers. $L_1$ has 256 units; $L_2$ has 128
% 	\item 10 output variables due to class labeling
% }

% \end{columns}

% \end{frame}
% %--- Next Frame ---%

% \begin{frame}[t]{More on that architecture}

% \begin{columns}
% \column{.25\textwidth}
% \centering
% \includegraphics[width = \textwidth, align = c]{Fig10-4}
% \column{.45\textwidth}
% \centering
% \includegraphics[width = \textwidth, align = c]{Eq10-10}
% \includegraphics[width = \textwidth, align = c]{Eq10-11}
% \InstructorList{
%   \item Superscript denotes layer
% 	\item $W_1$ denotes entire matrix of weight
% 	\item In this setting, size is $785 \times 256 = 200,960$ values
% 	\item 785 instead of 784 to involve intercept term (called bias in this literature)
% }

% \end{columns}

% \end{frame}
% %--- Next Frame ---%
% \begin{frame}[t]{Matrix version: First layer}
% \begin{columns}
% \column{.25\textwidth}
% \centering
% \includegraphics[width = \textwidth, align = c]{Fig10-4}

% \column{.65\textwidth}
% \centering
% \includegraphics[width = .6\textwidth, align = c]{Eq10-10}
% \begin{equation*}
% 	A^{(1)} = g( \mathbf{W}^{(1)}\cdot \mathbf{X})
% 	\qquad 
% 	\mathbf{X}^T = (1 \; X_1\; X_2\; \cdots \; X_p)
% \end{equation*}
% \begin{equation*}
% \end{equation*}
% \end{columns}
% \end{frame}
% %--- Next Frame ---%
% \begin{frame}[t]{Matrix version: Second layer}
% \begin{columns}
% \column{.25\textwidth}
% \centering
% \includegraphics[width = \textwidth, align = c]{Fig10-4}

% \column{.65\textwidth}
% \centering
% \includegraphics[width = .6\textwidth, align = c]{Eq10-11}
% \begin{equation*}
% 	A^{(2)} = g( \mathbf{W}^{(2)}\cdot \mathbf{A})
% 	\qquad 
% 	\mathbf{A}^T = (1 \; A^{(1)}_1\; A^{(1)}_2\; \cdots \; A^{(1)}_p)
% \end{equation*}
% \begin{equation*}
% \end{equation*}
% \end{columns}
% \end{frame}
% %--- Next Frame ---%
% \begin{frame}[t]{The last column for classification}

% \begin{columns}
% \column{.45\textwidth}
% \centering

% \includegraphics[width = \textwidth, align = c]{Fig10-4}
% \column{.45\textwidth}
% \centering
% \includegraphics[width = \textwidth, align = c]{Eq10-12}
% \includegraphics[width = \textwidth, align = c]{Eq10-13}
% \InstructorList{
%   \item Doing linear combo of inputs to get 10 different outputs
% 	\item Want this to act like a probability.
% 	\item Answer is to use the softmax activation function $f_m(X)$
% 	\item Values are non-negative and sum to 1
% 	\item Return class with highest probability
% }


% \end{columns}

% \end{frame}
% %--- Next Frame ---%
% \begin{frame}[t]{Matrix version: Last layer}
% 	\begin{columns}
% 	\column{.25\textwidth}
% 	\centering
% 	\includegraphics[width = \textwidth, align = c]{Fig10-4}
	
% 	\column{.65\textwidth}
% 	\centering
% 	\includegraphics[width = .6\textwidth, align = c]{Eq10-11}
% 	\begin{equation*}
% 		A^{(2)} = g( \mathbf{W}^{(2)}\cdot \mathbf{A})
% 		\qquad 
% 		\mathbf{A}^T = (1 \; A^{(1)}_1\; A^{(1)}_2\; \cdots \; A^{(1)}_p)
% 	\end{equation*}
% 	\begin{equation*}
% 	\end{equation*}
% 	\end{columns}
% 	\end{frame}
% %--- Next Frame ---%
% \begin{frame}[t]{MNIST}

% \begin{columns}
% \column{.45\textwidth}
% \centering
% \includegraphics[width = \textwidth, align = c]{Fig10-3}
% \column{.45\textwidth}
% \centering
% \InstructorList{
%   \item Goal: Build a model to classify images into their correct digit class
% 	\item Each image has $p=28 \cdot 28 = 784$ pixels
% 	\item Each pixel is grayscale value in [0,255]
% 	\item Data converted into column order
% 	\item Output represented by one-hot vector $Y = (Y_0,Y_1,\cdots,Y_9)$
% 	\item 60K training images, 10K test images
% }

% \end{columns}

% \end{frame}
% %--- Next Frame ---%
% \begin{frame}[t]{Results}

% \begin{columns}
% \column{.45\textwidth}
% \centering
% \includegraphics[width = .5\textwidth, align = c]{Fig10-3}
% \includegraphics[width = \textwidth, align = c]{Table10-1}
% \column{.45\textwidth}
% \centering
% \InstructorList{
%   \item Bottom two are previously covered
% 	\item Dramatic increase with neural nets
% 	\item Dropout regularization covered later..... maybe we won't get there
% 	\item Number of coefficients to train in $W_1$, $W_2$, and $B$ is 235,146; more than 33 times number needed for multinomial logistic regression
% 	\item $p \approx 4 n$, so need to avoid overfitting
% 	\item Regularization is important, but we won't get to that.

% }

% \end{columns}

% \end{frame}
% %--- Next Frame ---%



%======================


% \begin{frame}[t]{TL;DR}

% \begin{columns}
% \column{.45\textwidth}
% \centering
% \textbf{Feed Forward Neural Net}
% % \includegraphics[width = .5\textwidth, align = c]{Fig10-4}
% % \includegraphics[width = \textwidth, align = c]{Eq10-2}

% \begin{itemize}
% 	\item Combines input data using learned weights
% 	\item Linear combo of those to get output
% 	% \item Sometimes softmax to get probability of classification
% \end{itemize}

% \column{.45\textwidth}
% \centering
% % \textbf{CNN}
% % \includegraphics[width = \textwidth, align = c]{Fig10-8}

% % \begin{itemize}
% % 	\item Specialized NN
% % 	\item Gets next layer via
% % 	\begin{itemize}
% % 		\item Convolution layer
% % 		\item Pooling Layer
% % 		\item Fully connected layerw
% % 	\end{itemize}
% % \end{itemize}

% \end{columns}


% \end{frame}
% %--- Next Frame ---%

\begin{frame}[t]{Next time}

	\includegraphics[align = c, width = .4\textwidth]{Day01-Schedule-Test3}
	

\end{frame}
%--- Next Frame ---%


\end{document}
